% Provenance (never printed): sources/book2/BRIEFS.md, g11-09. The same assumption about incentives for voters, officials
% and agencies; rational ignorance; concentrated benefits and dispersed costs; rent seeking; capture; imperfect market
% against imperfect government. Real numbers: figdata/grade-11/09-public-choice-first-look/ (GAO/RCED-00-126 cited table
% and WDI population). Ledger: GAO 2000 (cost, benefit, farms, net loss), Motor Carrier Act 1980, Moore (trucking),
% Gelman, Silver and Edlin 2012. No published estimate of total lobbying or rent-seeking spending was read; none is
% printed. The second chart (Brenland's quota) is drawn as a diagram from the problem's numbers. Brenland is invented.
\chapter{Public Choice: A First Look}\label{ch:g11:public-choice-first-look}

In 1998 the American government's sugar programme, which kept the price of sugar in the United States above the world price by limiting imports, cost the country's
sugar users about 1.9 billion dollars. Spread over 276 million people, that was 6.89 dollars a person, too little for anyone to notice. Its benefit to sugar producers was
about 1.0 billion dollars; spread over the 8070 farms that grew sugar cane or sugar beet, that was up to 124 thousand dollars a farm. The programme had been in place for
decades, and it survived.
% figdata: sugar_key (pop 275.9 million, 8070 farms, 6.89, 123900); ledger: g11:gao-sugar-cost, g11:gao-sugar-benefit, g11:gao-sugar-farms

The earlier chapters asked how markets fail and how a government can correct them. This one asks how governments decide, using the same assumption the book uses for
buyers and firms: that voters, elected officials and agencies respond to their own costs and benefits. That assumption does not say they are selfish; it says they
respond to incentives. It explains why a programme like sugar's lasts.

\begin{openingchart}{The US sugar programme in 1998: its cost per US resident against its benefit per farm growing sugar cane or beet, in 1999 dollars, on a times-ten
scale. The benefit includes the processors' share, so the bar is an upper bound for a grower.}{US General Accounting Office (2000), report GAO/RCED-00-126; World Bank,
World Development Indicators (population)}
\begin{tikzpicture}
\begin{semilogyaxis}[width=0.62\linewidth, height=5.4cm, ybar, /pgf/bar width=34pt, xtick={0,1}, xticklabels={cost per person, benefit per farm},
  enlarge x limits=0.5, ymin=1, ymax=1000000, ytick={1,10,100,1000,10000,100000,1000000}, log ticks with fixed point, ylabel={dollars},
  y tick label style={/pgf/number format/1000 sep={\,}}, log origin=infty,
  nodes near coords={\pgfmathprintnumber[fixed, precision=2, 1000 sep={\,}]{\pgfplotspointmeta}}, point meta=rawy,
  every node near coord/.append style={font=\scriptsize}]
\addplot[fill=omDef!60, draw=omDef] table[x=x, y=dollars] {figdata/grade-11/09-public-choice-first-look/out/sugar.csv};
\end{semilogyaxis}
\end{tikzpicture}
\end{openingchart}

\begin{recall}
From the grades 6--9 volume: governments raise taxes and spend them on public services and transfers. From \cref{ch:g10:tax-incidence}: the
\omterm{def:g10:what-prices-do:dwl}{deadweight loss} of a tax. From \cref{ch:g11:oligopoly-and-games}: the \omterm{def:g11:oligopoly-and-games:prisoners-dilemma}{prisoner's dilemma}. From \cref{ch:g11:externalities}: the instruments that
correct an \omterm{def:g11:externalities:externality}{externality}, which a government must choose and set.
\end{recall}

\begin{definition}[Government failure]\label{def:g11:public-choice-first-look:government-failure}
\emph{Government failure}\index{government failure} occurs when a government's action lowers \omterm{def:g10:surplus-and-efficiency:total-surplus}{total surplus}, or fails to raise it as much as it could, because of the
incentives and information of those who decide.
\end{definition}

\section{The uninformed voter}\label{sec:g11:public-choice-first-look:voter}

A week before the American presidential election of 2008, the chance that one vote would decide it was estimated at about 1 in 60 million on average, and about 1 in
10 million in the closest states. If the outcome is worth a large sum to a voter, say 10\,000 dollars, the expected value of his vote changing it is
$10\,000/60\,000\,000$, less than a cent. Studying the candidates' programmes for many hours does not pay on such a stake. Most voters vote for other reasons (duty,
identity, the pleasure of taking part), but they have little reason to learn the details of the sugar programme.
% ledger: g11:gse2012

\begin{definition}[Rational ignorance]\label{def:g11:public-choice-first-look:rational-ignorance}
\emph{Rational ignorance}\index{rational ignorance} is a voter's choice not to learn about a policy, because the cost of learning is larger than the expected benefit,
which is small since one vote rarely decides anything.
\end{definition}

\section{Concentrated benefits, dispersed costs}\label{sec:g11:public-choice-first-look:interest-groups}

\begin{definition}[Interest group]\label{def:g11:public-choice-first-look:interest-group}
An \emph{interest group}\index{interest group} is an organised group that seeks to influence policy for its members' benefit: growers, an industry, a profession, a
union.
\end{definition}

\begin{definition}[Concentrated benefits and dispersed costs]\label{def:g11:public-choice-first-look:concentrated-benefits}
A policy has \emph{concentrated benefits and dispersed costs}\index{concentrated benefits and dispersed costs} when its gains go to a few, each of whom gains much,
and its costs fall on many, each of whom loses little.
\end{definition}

The few have both the motive and the means to organise: a large stake each, and few members to bring together. The many have neither. Each would gain from
organising against the policy, but each would rather others paid for it: organising is a \omterm{def:g11:public-goods-and-commons:private-good}{public good} for the losers, and they free ride.

\begin{model}{Concentrated benefits, dispersed costs}
\assumes Each group spends on influence up to what the policy is worth to it; organising a group costs more the larger and more scattered it is; voters are rationally
  ignorant of small costs.
\predicts Policies that give large gains to a few and small losses to many survive, even when the losses add up to more than the gains.
\tested Consistent with the US sugar programme, which cost users about 1.9 billion dollars in 1998 and brought producers about 1.0 billion; after the share
  that went to foreign suppliers and the losses from producing and consuming the wrong amounts, the country lost about 900 million dollars that year, and
  the programme lasted. The end of US trucking regulation in 1980, against the carriers and the drivers' union (record below), is a case the simple model
  does not predict. A test would need many programmes, with the ones that ended counted as well as the ones that lasted.
\breaks When an issue becomes salient, so that voters notice; when the costs also fall heavily on a few organised groups, which then lobby back.
\end{model}
% ledger: g11:gao-sugar-cost, g11:gao-sugar-benefit, g11:gao-sugar-net

\begin{record}[Sugar in the United States]
% ledger: g11:gao-sugar-cost, g11:gao-sugar-benefit, g11:gao-sugar-farms, g11:gao-sugar-net; figdata: sugar_key
A government audit of the programme estimated, in 1999 dollars, that it cost American users of sugar 1.5 billion dollars in 1996 and 1.9 billion in 1998, and brought
domestic producers 800 million and 1.0 billion. About 70\,\% of the producers' gain went to sugar beet growers and processors, 30\,\% to sugar cane. Of the net loss to
the country, about 300 million dollars in 1996 and 500 million in 1998 was surplus destroyed, and about 400 million each year was paid to foreign producers allowed to
sell sugar into the American market at the high domestic price. The audit's 1997 census count was 973 cane farms and 7097 beet farms.
\end{record}

The diagram below shows the same accounts for an invented country, Brenland, which limits sugar imports by a quota. The rise in the domestic price transfers surplus from
buyers to domestic growers (area $a$) and to the holders of import licences (area $c$); areas $b$ and $d$ are destroyed, as in any restriction of trade.

\begin{omfigure}
\begin{tikzpicture}[econ, x=0.0105cm, y=0.62cm]
\fill[omDef!20] (0,2) -- (100,2) -- (300,4) -- (0,4) -- cycle; \node[font=\small] at (120,3) {$a$};
\fill[omThm!25] (100,2) -- (300,2) -- (300,4) -- cycle; \node[font=\small] at (240,2.6) {$b$};
\fill[omProp!25] (300,2) rectangle (500,4); \node[font=\small] at (400,3) {$c$};
\fill[omThm!25] (500,2) -- (700,2) -- (500,4) -- cycle; \node[font=\small] at (560,2.6) {$d$};
\draw[econ axes] (0,9.5) node[left, font=\small] {dalens a kg} |- (950,0) node[below left=12pt and 0pt, font=\small] {thousand tonnes};
\draw[demand] (0,9) -- (900,0) node[above right, font=\scriptsize] {demand};
\draw[supply] (0,1) -- (800,9) node[right, font=\scriptsize] {domestic supply};
\draw[guide] (0,2) -- (850,2) node[right, font=\scriptsize] {world price};
\draw[guide] (0,4) -- (500,4);
\foreach \x in {100,300,500,700} \draw[guide] (\x,0) -- (\x,2);
\draw[guide] (300,2) -- (300,4); \draw[guide] (500,2) -- (500,4);
\foreach \x in {100,300,500,700} \node[below, font=\footnotesize] at (\x,0) {\x};
\foreach \y in {2,4} \draw (0,\y) -- (-8,\y) node[left, font=\footnotesize] {\y};
\end{tikzpicture}

{\small Brenland's sugar market (the weekend problem). At the world price of 2 dalens a kilogram, growers produce 100 thousand tonnes and buyers want 700. A quota of 200
thousand tonnes of imports raises the domestic price to 4. Growers gain $a$ = 400 million dalens; licence holders $c$ = 400; buyers lose $a + b + c + d$ = 1200; $b + d$
= 400 is the \omterm{def:g10:what-prices-do:dwl}{deadweight loss}.}
\end{omfigure}

\section{Rent seeking}\label{sec:g11:public-choice-first-look:rent-seeking}

\begin{definition}[Rent seeking]\label{def:g11:public-choice-first-look:rent-seeking}
\emph{Rent seeking}\index{rent seeking} is spending resources to obtain a transfer, a protected position or a licence from the government, rather than to produce
something of value.
\end{definition}

The cost of a programme is not only its \omterm{def:g10:what-prices-do:dwl}{deadweight loss}. Firms competing for a licence or a quota spend on lobbyists, lawyers and campaigns, and that spending is lost
to everyone. Suppose two firms compete for an import licence worth $V$, and each wins with a probability equal to its share of the total spent. If firm 1 spends $x_1$
and firm 2 spends $x_2$, firm 1 expects $V x_1/(x_1 + x_2) - x_1$. Its best spending satisfies $V x_2/(x_1 + x_2)^2 = 1$; in the symmetric equilibrium $x_1 = x_2 = x$,
this gives $V/(4x) = 1$, so each spends $V/4$ and together they spend half the prize. With more competitors, more of the prize is spent getting it.

\begin{model}{Rent seeking}
\assumes A transfer goes to whoever spends most on obtaining it, or with a chance that rises with spending; contestants are free to compete.
\predicts Spending on the contest rises towards the value of the prize; the \omterm{def:g11:externalities:private-cost}{social cost} of a transfer is its \omterm{def:g10:what-prices-do:dwl}{deadweight loss} plus the resources spent seeking it.
\tested Hard to measure, since much of the spending is not recorded as such. Licences that could be sold had prices: in the 1970s the right to carry certain goods on
  certain routes by truck sold for hundreds of thousands of dollars, a measure of the rent the regulation created.
\breaks When the rules are open to all bidders (an auction of the licence turns the rent into government revenue instead of lobbying), or when rents are small.
\end{model}
% ledger: g11:moore-trucking

\section{Capture}\label{sec:g11:public-choice-first-look:capture}

\begin{definition}[Regulatory capture]\label{def:g11:public-choice-first-look:capture}
\emph{Regulatory capture}\index{regulatory capture} occurs when a regulator comes to serve the interests of the industry it regulates rather than those of the public it
was set up to protect.
\end{definition}

A regulator deals every day with the firms it oversees, depends on them for information, and may hire from them or move to them. The firms are few and attentive; the
public is many and rationally ignorant. Capture follows from the same incentives as the sugar programme.

\begin{record}[Capture: trucking in the United States until 1980]
% ledger: g11:moore-trucking, g11:mca-1980
For decades a federal commission set American trucking rates and decided who could carry what on which route. Rights to routes became valuable, selling for hundreds of
thousands of dollars by the 1970s, and unionised drivers earned about 50\,\% more than comparable workers. Both the truckers' association and the drivers' union strongly
opposed ending the regulation. The Motor Carrier Act, signed on 1 July 1980, opened entry and pricing. Truckload rates fell about 25\,\% in real terms between 1977 and
1982, the value of route rights fell close to zero, and by 1990 there were more than 40\,000 licensed carriers, over twice as many as in 1980.
\end{record}

\section{Imperfect markets, imperfect governments}\label{sec:g11:public-choice-first-look:comparison}

The earlier chapters showed real market failures with measured costs: pollution that markets do not price, fisheries emptied, insurance pools that unravel. This
chapter shows that the cure is chosen by people with their own incentives. The comparison that matters is never a failing market against a perfect government, or a
failing government against a perfect market, but one imperfect arrangement against another, each judged by its evidence. A \omterm{def:g11:externalities:pigouvian-tax}{Pigouvian tax} that a lobby can lower, a
quota that its holders can capture, a regulator that the regulated can win over: each is a case to be measured, not assumed.

\section{Exercises}

\begin{exercise}[$\star$]\label{exo:g11:public-choice-first-look:1}
A programme costs consumers 600 million dalens a year; Brenland has 3 million households. What does it cost each household?
\end{exercise}

\begin{exercise}[$\star$]\label{exo:g11:public-choice-first-look:2}
The same programme gives 300 million dalens a year to 1500 producers. What does each gain? How many times the cost to a household is that?
\end{exercise}

\begin{exercise}[$\star$]\label{exo:g11:public-choice-first-look:3}
From the opening chart, how many times larger was the benefit per farm than the cost per person? Why is the benefit per farm an upper bound for a grower?
\end{exercise}

\begin{exercise}[$\star$]\label{exo:g11:public-choice-first-look:4}
For each policy, who has concentrated benefits and who has dispersed costs? A tariff on imported steel; a licence required to work as a hairdresser; a subsidy to
a country's few aircraft makers.
\end{exercise}

\begin{exercise}[$\star$]\label{exo:g11:public-choice-first-look:5}
A voter's vote decides an election with probability 1 in 10 million, and the better candidate's win is worth 50\,000 dalens to her. What is the expected value of her
vote changing the outcome? Is it worth 10 hours of study at 20 dalens an hour?
\end{exercise}

\begin{exercise}[$\star\star$]\label{exo:g11:public-choice-first-look:6}
Two firms compete for a licence worth 800 thousand dalens; each wins with a probability equal to its share of the total spent. Find each firm's spending and the total in
equilibrium.
\end{exercise}

\begin{exercise}[$\star\star$]\label{exo:g11:public-choice-first-look:7}
With $n$ identical firms in the contest of \cref{exo:g11:public-choice-first-look:6}, each spends $V(n-1)/n^2$. Check this for $n = 2$, and compute the total spent for
$n = 4$ and $n = 10$. What happens as $n$ grows?
\end{exercise}

\begin{exercise}[$\star\star$]\label{exo:g11:public-choice-first-look:8}
Brenland lowers its sugar quota of the diagram from 200 to 100 thousand tonnes. Find the new domestic price, the growers' gain and the \omterm{def:g10:what-prices-do:dwl}{deadweight loss}.
\end{exercise}

\begin{exercise}[$\star\star$]\label{exo:g11:public-choice-first-look:9}
Learning how a programme affects her costs a voter 200 dalens of time. The programme costs her 30 dalens a year, and her vote is decisive with probability 1 in a million.
Should she learn? What would change her answer?
\end{exercise}

\begin{exercise}[$\star\star$]\label{exo:g11:public-choice-first-look:10}
In the diagram of Brenland's sugar market, read the areas $a$, $b$, $c$, $d$, and check that the buyers' loss equals their sum.
\end{exercise}

\begin{exercise}[$\star\star$]\label{exo:g11:public-choice-first-look:11}
A newspaper campaign makes the sugar programme of the opening chart known to most voters. Which assumption of the interest-group model does this change, and what does the
model then predict?
\end{exercise}

\begin{exercise}[$\star\star\star$]\label{exo:g11:public-choice-first-look:12}
What evidence would contradict the interest-group model? Use the trucking record: who gained from the regulation, who opposed its end, and why its end is a puzzle
for the model. Suggest what the model leaves out.
\end{exercise}

\begin{exercise}[$\star\star\star$]\label{exo:g11:public-choice-first-look:13}
A regulator's official earns 60 thousand dalens a year and can expect a job in the regulated industry at 150 thousand if the industry finds him helpful, with probability
$\pi$. A ruling in the industry's favour raises $\pi$ from 0.1 to 0.4. What is that ruling worth to him, for a job that lasts five years? What rules could weaken this
incentive?
\end{exercise}

\begin{exercise}[$\star\star\star$]\label{exo:g11:public-choice-first-look:14}
A factory's pollution causes a \omterm{def:g10:what-prices-do:dwl}{deadweight loss} of 100 million dalens a year. A tax would remove it, but the lobbying it provokes is expected to set the tax at twice the
right level, causing a \omterm{def:g10:what-prices-do:dwl}{deadweight loss} of 100 million of its own, and the contest costs 30 million. Compare the outcomes. What does the comparison depend on?
\end{exercise}

\begin{exercise}[$\star\star\star$]\label{exo:g11:public-choice-first-look:15}
Three legislators vote on two projects, each costing 30 to each of them in taxes. Project X brings benefits of 80, 10 and 10 to legislators 1, 2 and 3; project Y brings
10, 80 and 10. Does each project pass a majority vote alone? What happens if legislators 1 and 2 agree to vote for both? Is the outcome efficient?
\end{exercise}

\section{Problem: Brenland's Sugar Quota}

\begin{problem}[{Weekend problem --- who gains, who pays and what it costs}]
\label{pb:g11:public-choice-first-look:1}
In Brenland, demand for sugar is $Q_d = 900 - 100p$ and domestic supply $Q_s = 100p - 100$ (thousand tonnes a year; $p$ in dalens a kilogram). The world price is 2.
Brenland has 200 sugar growers and 3 million households. A thousand tonnes at one dalen a kilogram is one million dalens.

\medskip
\noindent\textbf{Part I --- Free trade.}
\begin{enumerate}
\item Find consumption, domestic production and imports at the world price.
\item Compute the growers' \omterm{def:g10:surplus-and-efficiency:producer-surplus}{producer surplus}.
\item Compute the buyers' \omterm{def:g10:surplus-and-efficiency:consumer-surplus}{consumer surplus}.
\end{enumerate}

\medskip
\noindent\textbf{Part II --- A quota.}
\begin{enumerate}[resume]
\item Imports are limited to 200 thousand tonnes. Find the domestic price.
\item Find consumption and domestic production.
\item Compute the gain of the growers (area $a$).
\item Compute the quota rent of the licence holders (area $c$).
\item Compute the buyers' loss.
\item Compute the \omterm{def:g10:what-prices-do:dwl}{deadweight loss} ($b + d$) and check that the accounts add up.
\item Which areas are lost to Brenland if the licences go to foreign exporters?
\end{enumerate}

\medskip
\noindent\textbf{Part III --- Who organises.}
\begin{enumerate}[resume]
\item Compute the gain per grower and the cost per household.
\item How many times the cost per household is the gain per grower?
\item The growers' association spends 40 million dalens a year lobbying for the quota. What share of the growers' gain is that?
\item Is it worth a household's while to spend one hour, worth 20 dalens, writing to its representative? Explain with the free-rider problem.
\item Firms compete for the import licences in a contest like the chapter's, with two equal contestants. How much do they spend in all?
\end{enumerate}

\medskip
\noindent\textbf{Part IV --- The full cost.}
\begin{enumerate}[resume]
\item Add the \omterm{def:g10:what-prices-do:dwl}{deadweight loss}, the growers' lobbying and the licence contest of questions 13 and 15. What is the full yearly cost of the quota to Brenland?
\item What share of the buyers' loss reaches the growers as a gain, net of their lobbying?
\item If the licences were auctioned, what would change in your answer to question 16?
\item Suggest one change that would make the quota's cost salient to voters.
\item State the result of the problem: the gain per grower against the cost per household, and the share of the growers' gain spent on lobbying.
\end{enumerate}
\end{problem}
